% Generated draft: requires mathematical and source review.
\documentclass[11pt]{article}
\usepackage[margin=25mm]{geometry}
\usepackage{fontspec}
\setmainfont{Latin Modern Roman}
\usepackage{amsmath,amssymb,mathtools,unicode-math}
\setmathfont{Latin Modern Math}
\usepackage{enumitem,needspace,xurl,hyperref}
\setlength{\parindent}{0pt}
\setlength{\parskip}{5pt}
\setlength{\emergencystretch}{3em}
\newcommand{\N}{\mathbb{N}}
\newcommand{\Z}{\mathbb{Z}}
\newcommand{\Q}{\mathbb{Q}}
\newcommand{\R}{\mathbb{R}}
\newcommand{\C}{\mathbb{C}}
\newcommand{\F}{\mathbb{F}}
\newcommand{\A}{\mathbb{A}}
\newcommand{\PP}{\mathbb{P}}
\newcommand{\E}{\mathbb{E}}
\newcommand{\eps}{\varepsilon}
\newcommand{\dd}{\,\mathrm{d}}
\begin{document}
\section*{UnsolvedMath ID: 800971}
\textbf{Source problem code:} catalogue entry
\subsection{Definitions and mathematical statement}
A Lorentzian manifold is a smooth manifold \(M\) equipped with a smooth nondegenerate symmetric metric \(g\) of index one, up to reversing the overall sign convention. Here \emph{closed manifold} means compact and without boundary. No orientability, time-orientability, or geodesic-completeness hypothesis is imposed in the question as supplied.

Under the periodic affine interpretation, a \emph{closed geodesic} is a nonconstant smooth map \(\gamma:\R\to M\) for which there exists \(T>0\) such that
\[
\nabla_{\dot\gamma}\dot\gamma=0,
\qquad
\gamma(t+T)=\gamma(t)\quad\text{for every }t\in\R,
\]
where \(\nabla\) is the Levi--Civita connection of \(g\). With this interpretation, the question is:
\[
\text{Does every closed Lorentzian manifold }(M,g)
\text{ admit such a nonconstant periodic geodesic?}
\]
The source does not explicitly specify whether closure is required for an affine parametrization or only for an unparametrized geodesic. The displayed formulation uses the former interpretation; this distinction requires review against the source. The requested geodesic is not restricted to any causal type.

The supplied background states that a closed orientable surface admitting a metric of signature \((+,-)\) is homeomorphic to the torus. It attributes to [Suhr2013a] the assertion that every Lorentzian two-dimensional torus has at least two simple closed geodesics, at least one of which is definite, meaning timelike or spacelike, and states that explicit examples show this assertion to be optimal. A simple closed curve is one whose parametrization from a circle is an embedding; a definite geodesic has \(g(\dot\gamma,\dot\gamma)\) everywhere nonzero with constant sign. These background claims are retained as supplied attribution, without independent verification. The source's statement that, to its authors' knowledge, nothing was known in dimensions \(\geq 3\) is a historical knowledge qualification, not a verified current literature assessment.
\subsection{Short English statement}
Does every compact Lorentzian manifold without boundary contain a nonconstant closed geodesic? The displayed formulation requires a periodic affine parametrization; the intended meaning of closure needs source review. The registry labels the problem open, and the supplied background reports a stronger existence result for Lorentzian two-dimensional tori.
\subsection{Sources}
\begin{itemize}
\item Local UnsolvedMath registry, permanent ID \(800971\), exact source problem code \texttt{catalogue entry}; collection: AMR Open Problem Lists. The record attributes the problem to Burns and Matveev, gives proposed year 2019, and supplies status \emph{open}. These are registry metadata, not an independently verified status assessment.
\item Burns and Matveev, \emph{Open Problems and Questions about Geodesics} (2019), Question 52, as identified in the supplied record: \url{https://arxiv.org/abs/1308.5417}. The record reports access on 2026-10-16 and extraction by \texttt{arxiv\_latex}; that access is supplied provenance and is not claimed here as an independent source visit.
\item [Suhr2013a], cited in the supplied statement for the Lorentzian two-torus result. No full bibliographic entry or direct link for this reference is supplied, so its exact identity and hypotheses remain to be checked.
\end{itemize}
\end{document}
